%%%%%%%%%%%%%%%%%%%%%%%%%%%%%%%%%%%%%%%%%%%%%%%%%%%%%%%%%%%%%%%%%%%%%%
%%                       IAEW LaTeX Template                        %%
%%%%%%%%%%%%%%%%%%%%%%%%%%%%%%%%%%%%%%%%%%%%%%%%%%%%%%%%%%%%%%%%%%%%%%

\chapter{Diskussion}
\label{ch:discussion}
% ===============================================================
% STAND 24.09.2026, zweite Fassung des Tages. Sie ersetzt die Gliederung
% mit acht Abschnitten vollstaendig. Anlass ist die Vorgabe des
% Verfassers, die Diskussion an der Logik des Gesamtkonzepts
% auszurichten statt an den Abschnitten des Kapitels 4.
%
% KAPITELTITEL: nur noch "Diskussion" statt "Bewertung und Diskussion".
%
% STRUKTUR, Vorgabe des Verfassers im Wortlaut nachgebildet:
%   5.1 Anforderungen und ihre Umsetzung im Produkt
%   5.2 Modellierung und ihre Annahmen
%   5.3 Wirtschaftlichkeit und Wirksamkeit
%   5.4 Batteriespeicher im Engpassmanagement und am Markt
%   5.5 Umsetzung, systemweite Ausrollung und Forschungsbedarf
%
% ZAHLEN: Vorgabe des Verfassers vom 24.09.2026. Die Diskussion nennt
% fast keine Zahlen, sondern die Erkenntnisse, Schluesse und
% Analyseergebnisse, die aus den Zahlen folgen. Stehen bleiben allein
% Produktgroessen wie die Reaktionszeitklasse und Paragrafen. Jede Zahl,
% die doch noetig wird, ist beim Ausformulieren einzeln zu begruenden.
%
% EINLEITUNG: beschreibt genau einmal den Aufbau des Kapitels. Die
% Bestandsaufnahme ist auf die tragenden Erkenntnisse verdichtet und
% traegt keine Zahlen mehr.
%
% WOHER DIE STICHPUNKTE STAMMEN. Die Gliederung mit acht Abschnitten vom
% 24.09.2026 vormittags steht im Git unter dem Commit e8c4012, ebenso die
% ausformulierte Kapiteleinleitung jener Fassung. Das Stichpunktgeruest
% vom 18.08.2026 steht am Ende dieser Datei auskommentiert und traegt den
% Belegstand jedes Punktes mit den Kennungen B1 bis G4.
%
% NEU AUFGENOMMEN aus Kapitel 3, Vorgabe des Verfassers vom 24.09.2026,
% die Annahmen und offenen Punkte, die dort benannt und bisher nicht
% diskutiert sind:
%   Nachlaufpflicht offen (3.1.2), Hoehe der Abrufverguetung offen
%   (3.1.2), Mindestgroesse und Zusammenfassung je Netzknoten (3.1.2),
%   pay-as-bid und Zuschlag als Verbund beider Richtungen (3.1.2), kein
%   Verguetungsrahmen vorhanden (3.1.2), tagesweise Loesung ohne
%   Ladezustandsuebergang (3.2.1), Netz, Pool und Gegenseite nicht
%   abgebildet (3.2.1), Rueckwirkung der Speicher auf die Preise ausser
%   Betracht (3.2.1), Ausschluss der Intraday-Auktionen (3.2.3), kein
%   mehrfacher Handel derselben Viertelstunde (3.2.3), Validierung
%   schliesst allein eine Unterschaetzung aus (3.3.2).
%
% STILVORGABEN wie im ganzen Dokument, siehe CLAUDE.md Abschnitt 4.
% Die fett gesetzten Absatzkoepfe sind Geruest und fallen mit der
% Ausformulierung weg.
%
% BEGRIFFE, die dieses Kapitel voraussetzt, mit ihrer Definitionsstelle:
%   kurativer Reservierungspreis .......... 1.2
%   kurative Reservierung, kurative Bindung  2.1.2
%   Referenzfahrplan, Referenzerloes ...... 2.1.2
%   Indifferenzprinzip .................... 2.3.1
%   Arbitrage ............................. 2.2.1
%   Anforderungskatalog A1 bis A8 ......... 3.1.1
%   Bezugsanlage .......................... 3.2.2
%   Abrufdauer, Bindungsdauer, Reaktionszeit  2.1.2
% NEU EINZUFUEHREN in diesem Kapitel, mit je einem Satz:
%   kapazitaetsbasierter Redispatch ....... 5.5, Beleg noch zu pruefen
% ===============================================================

\noindent\textbf{Absatz 1, was vorliegt und was daraus einzuordnen ist.}
\begin{itemize}
    \item Nach der Ableitung der Anforderungen, dem Entwurf der kurativen Reservierung und der Bestimmung des kurativen Reservierungspreises liegen Erkenntnisse vor, die in das Gesamtkonzept eines kurativen Marktmechanismus einzuordnen sind
    \item Der kurative Reservierungspreis entsteht für jede Stunde neu und streut über das Jahr weit, sodass eine einzelne Zahl ihn nicht wiedergibt
    \item Er richtet sich vor allem gegen die Vorhaltung von Regelleistung und nur in wenigen Stunden gegen den Handel
    \item Er ist im Winter am günstigsten, und dort liegt zugleich der größte Engpassmanagementbedarf
    \item Er ist die Untergrenze eines Gebots, denn die Wahrscheinlichkeit eines Abrufs, die Sanktion und der Aufwand des Anbietens fehlen ihm
\end{itemize}

\noindent\textbf{Absatz 2, wie dieses Kapitel vorgeht.}
% Dieser Absatz beschreibt den Aufbau GENAU EINMAL. Kein weiterer
% Abschnitt wiederholt ihn.
\begin{itemize}
    \item Zuerst werden die Anforderungen aus Abschnitt~\ref{sec:requirements_catalogue} durchgegangen und mit ihrer jeweiligen Umsetzung im Produkt verglichen
    \item Anschließend wird geprüft, ob die Annahmen der Modellierung ausreichen und wie einzelne Annahmen auf einen möglichen kurativen Reservierungspreis wirken
    \item Darauf folgt die Bewertung der Wirtschaftlichkeit anhand des Preises und der Wirksamkeit im Zusammenhang mit der Einsparung von Redispatch
    \item Anschließend wird der Zusammenhang zwischen \ac{BESS} und den kurzfristigen Märkten behandelt, dazu ihre künftige Rolle im Strommarkt und bei den Systemdienstleistungen
    \item Zuletzt werden der Aufwand der Umsetzung und eine systemweite Ausrollung der kurativen Systemführung abgewogen, woraus der weitere Forschungsbedarf folgt
\end{itemize}

% ---------------------------------------------------------------
\section{Anforderungen und ihre Umsetzung im Produkt}
\label{sec:requirements_assessment}
% Ziel rund 2,5 Seiten, vier Absaetze.
% Fragen: Welche der acht Anforderungen loest der Entwurf ein? Welche
% beantwortet erst das Ergebnis? Welche bleibt offen? Was tritt hinzu,
% das der Katalog nicht fuehrt?
% TRAGENDER BEFUND: A7 Vertraeglichkeit.

\noindent\textbf{Absatz 1, vier Anforderungen entscheidet der Zuschnitt des Produkts.}
\begin{itemize}
    \item Übergang, in Abschnitt~\ref{sec:requirements_catalogue} werden acht Anforderungen abgeleitet, an denen die kurative Reservierung zu messen ist
    \item Die Reaktionszeitklassen differenzieren das Produkt nach dem betroffenen Betriebsmittel, womit A1 eingelöst ist
    \item Die Vergütung setzt an der vorgehaltenen Leistung und an der Bindungsdauer an, und ein Abruf wird gesondert vergütet, womit A4 eingelöst ist
    \item Die Präqualifikation prüft technische Eigenschaften und nicht die Art der Anlage, womit A5 dem Zugang nach eingelöst ist
    \item Die Beschaffung ist netzknotenscharf entworfen, womit A3 im Produkt eingelöst und im Ergebnis nicht geprüft ist
    \item Über diese vier Anforderungen entscheidet damit der Zuschnitt des Produkts und nicht der ermittelte Preis
\end{itemize}

\noindent\textbf{Absatz 2, die Teilnahme nach A6 lohnt sich.}
\begin{itemize}
    \item Übergang, ob der kurative Marktmechanismus Anbieter findet, entscheidet A6
    \item Unter den Preisen der vollständigen Bindung steht der Betreiber besser als ohne die Reservierung
    \item Die Teilnahme lohnt sich damit, sofern der \ac{ÜNB} den geforderten Preis zahlt
    % EIGENSTAENDIGE ABLEITUNG, vom Verfasser zu bestaetigen: die Summe der
    % stuendlichen Indifferenzpreise uebersteigt die entgangene Vermarktung,
    % weil jede Stunde gegen ihre eigene beste Verwendung bepreist wird,
    % waehrend die Verwendungen um dieselbe Leistung konkurrieren.
    \item Der Grund liegt darin, dass jede Stunde gegen ihre eigene beste Verwendung bepreist wird, während die Verwendungen um dieselbe Leistung konkurrieren
    \item Der Befund gilt für die Bezugsanlage, denn der Preis folgt ihren Erlösmöglichkeiten
    \item Für andere Technologien und andere Netzknoten ist die Teilnahme gesondert zu prüfen
\end{itemize}

\noindent\textbf{Absatz 3, die Verträglichkeit nach A7 ist der kritische Punkt.}
\begin{itemize}
    \item Übergang, die schärfste der acht Anforderungen ist A7
    \item A7 verlangt, dass der kurative Marktmechanismus neben den bestehenden Märkten besteht, ohne einen von ihnen zu verdrängen
    \item Unter vollständiger Bindung verdrängt die kurative Reservierung die gesamte übrige Vermarktung der Bezugsanlage
    \item Am stärksten trifft das die Vorhaltung von \ac{aFRR}-Leistung, auf die der größte Teil des Erlöses entfällt
    \item Würde eine große Zahl geeigneter Anlagen kurativ gebunden, verlöre der Markt für \ac{aFRR}-Leistung seine Anbieter
    \item A7 ist damit nicht über den Zuschnitt des Produkts zu erfüllen, sondern allein über eine Begrenzung der beschafften Menge
\end{itemize}

\noindent\textbf{Absatz 4, was offen bleibt und was der Katalog nicht führt.}
\begin{itemize}
    \item Übergang, zwei Anforderungen lassen sich aus dem Entwurf nicht beantworten
    \item A2 verlangt eine Sanktion, die das Produkt benennt, deren Höhe es aber nicht bemisst
    \item A8 betrifft die internen Abläufe des \ac{ÜNB} und ist von außen nicht zu beurteilen
    \item Hinzu tritt eine Konkurrenz, die der Katalog nicht führt, denn die Präqualifikation der Regelleistung reserviert bereits ein Ladezustandsband \cite{ubertragungsnetzbetreiber_deutschland_praqualifikationsverfahren_2024}
    \item A5 sichert zudem den Zugang und nicht den Wettbewerb, denn an einem Netzknoten mit wenigen wirksamen Anlagen liegt das Gebot über dem ermittelten Preis
    \item Die Zahl der Anbieter je Netzknoten entscheidet damit über die Lage des Preises
\end{itemize}

% ---------------------------------------------------------------
\section{Modellierung und ihre Annahmen}
\label{sec:model_discussion}
% Ziel rund 3 Seiten, fuenf Absaetze.
% Fragen: Welche Annahme traegt das Ergebnis? Wie wirkt die vollstaendige
% Preiskenntnis? Was misst die Sensitivitaet zur Handelsspanne nicht?
% Was laesst das Suchverfahren offen? Wie weit traegt die Perspektive und
% was belegt die Validierung?

\noindent\textbf{Absatz 1, die Modellierung der \ac{aFRR} trägt das Ergebnis.}
\begin{itemize}
    \item Übergang, der kurative Reservierungspreis folgt aus dem Optimierungsmodell, sodass dessen Annahmen zu prüfen sind
    \item Unter den geprüften Annahmen trägt die Modellierung der \ac{aFRR} das Ergebnis am stärksten
    \item Der Abrechnungspreis der Kapazitätsauktion bestimmt dabei die teuren Stunden, der Liefermodus die typische Stunde
    \item Die anteilige Lieferung des Basisfalls ist die belastbare Annahme, denn eine einzelne Anlage zieht den deutschlandweiten Abruf nicht auf sich
    \item Ein freier Liefermodus ließe das Optimierungsmodell allein die ertragreichsten Abrufstunden wählen, was ein Betreiber ohne Kenntnis des Abrufs nicht kann
    \item Sicher entscheiden ließe sich die Frage allein an den Gebotslisten der Auktionen, die anonymisiert sind
\end{itemize}

\noindent\textbf{Absatz 2, die vollständige Preiskenntnis verzerrt in beide Richtungen.}
\begin{itemize}
    \item Übergang, die stärkste Annahme des Optimierungsmodells ist die vollständige Kenntnis der Preise
    \item Sie überschätzt die Opportunität, weil ein Betreiber die Preise des folgenden Tages nicht kennt
    \item Weber räumt dieselbe Verzerrung für sein eigenes Verfahren ein, und der Unterschied zwischen Erlöspotenzial und erzieltem Erlös ist unabhängig davon belegt \cite{weber_gutachten_2015, mahgoub_wie_2025}
    \item Zugleich unterschlägt die vollständige Preiskenntnis den Optionswert der Bindung, denn bei bekannter Preisreihe ist er null
    \item Die Annahme ist gleichwohl bewusst gewählt, weil eine Preisprognose einen Fehler hinzufügte, der sich von der Wirkung der Reservierung nicht trennen ließe
    \item Beide Verzerrungen laufen gegeneinander, sodass der ermittelte Preis als opportunitätskostenbasierter Wert zu lesen ist
\end{itemize}

\noindent\textbf{Absatz 3, was die Sensitivität zur Handelsspanne nicht misst.}
% VORGEMERKT am 24.09.2026, Vorgabe des Verfassers: zur Wahrheit gehoert,
% dass die Sensitivitaet in 4.6.2 nicht misst, was sie messen soll.
\begin{itemize}
    \item Übergang, auch eine Sensitivität kann die Frage verfehlen, die sie beantworten soll
    \item Der Faktor auf die Handelsspanne streckt jede Abweichung des Intradaypreises von seinem Tagesmittel
    \item Er hebt damit gerade die Stunden, die ohnehin weit vom Tagesmittel liegen
    \item Der Vorteil des mehrfachen Handels entsteht dagegen innerhalb desselben Lieferzeitpunkts
    \item Beide Mengen von Stunden fallen nicht zusammen, sodass die Sensitivität allein die Breite der Handelsspanne misst
    \item Für eine Aussage über den Vorteil des mehrfachen Handels wären die Orderbücher des \ac{IDC} nötig
\end{itemize}

\noindent\textbf{Absatz 4, was das Suchverfahren offenlässt.}
% NICHT ALS FEHLER FORMULIEREN, sondern als Eigenschaft des Verfahrens.
\begin{itemize}
    \item Übergang, auch das Verfahren der Preissuche setzt dem Ergebnis eine Grenze
    \item Die Preise der zweiten Iteration sind ein koordinatenweises und kein globales Minimum
    \item Keine einzelne Stunde lässt sich senken, ohne dass die vollständige Bindung bricht
    \item Die Lage des Minimums hängt gleichwohl an der Reihenfolge des Abstiegs
    \item Der ausgewiesene Wert ist damit die vorsichtigere Seite einer gemessenen Spanne
\end{itemize}

\noindent\textbf{Absatz 5, die Grenzen der Perspektive und der Validierung.}
\begin{itemize}
    \item Übergang, die letzte Grenze ist die Perspektive der Untersuchung
    \item Das Optimierungsmodell rechnet eine einzelne Anlage und bildet weder das Netz noch die Zusammenfassung mehrerer Anlagen noch die Rückwirkung der Speicher auf die Preise ab
    \item Es löst jeden Liefertag für sich, sodass kein Ladezustand in den Folgetag übergeht
    \item Die Wahrscheinlichkeit eines Abrufs und die Sanktion bleiben außer Betracht, sodass der ermittelte Preis allein den opportunitätskostenbasierten Teil eines Gebots bildet
    \item Die Validierung stellt ein Modell mit vollständiger Preiskenntnis einem Verfahren mit rollierender Vorausschau gegenüber und einen engeren Marktumfang einem weiteren
    \item Sie schließt damit eine Unterschätzung der Erlösmöglichkeiten aus und belegt nicht die Höhe einer einzelnen Stunde
\end{itemize}

% ---------------------------------------------------------------
\section{Wirtschaftlichkeit und Wirksamkeit}
\label{sec:economics_effectiveness}
% Ziel rund 3 Seiten, vier Absaetze.
% Fragen: Was sagt der Preis wirtschaftlich aus? Gegen welchen Markt
% entsteht er? Wie verhaelt er sich zum Engpassmanagement? Welcher
% Zuschnitt und welcher Verguetungsrahmen folgen daraus?

\noindent\textbf{Absatz 1, was der ermittelte Preis wirtschaftlich aussagt.}
\begin{itemize}
    \item Übergang, der ermittelte Preis ist zuerst eine wirtschaftliche Größe
    \item Er misst die entgangene Marktopportunität derselben Stunde und folgt damit dem Indifferenzprinzip aus Abschnitt~\ref{sec:redispatch_compensation}
    \item Die Wahrscheinlichkeit eines Abrufs, die Sanktion und der Aufwand des Anbietens fehlen ihm
    \item Ein Betreiber böte deshalb über dem ermittelten Preis, und zwar um einen Aufschlag, den diese Arbeit nicht beziffert
    \item Die Hälfte der Stunden ist günstig zu reservieren, während einzelne Stunden den Durchschnitt tragen
    \item Gerade diese teuren Stunden lassen sich am schlechtesten vorhersagen, denn sie hängen an einzelnen Preisspannen des Energiemarkts
\end{itemize}

\noindent\textbf{Absatz 2, der Preis entsteht gegen die Regelleistung.}
\begin{itemize}
    \item Übergang, gegen welche Vermarktung der Preis entsteht, entscheidet über seine Höhe
    \item Der Leistungspreis der \ac{aFRR} entfällt mit der reservierten Stunde vollständig, weil er an eine Zeitscheibe von vier Stunden gebunden ist
    \item Die Arbitrage an den Energiemärkten lässt sich dagegen zwischen den Stunden eines Tages verschieben, solange der Tag eine Staffel gleichwertiger Spannen trägt
    \item Der Energiemarkt setzt den Preis deshalb nur an den wenigen Tagen, an denen diese Staffel einbricht
    \item Derselbe Mechanismus erklärt die wechselnde Verdrängungsreihenfolge, die offenen Stunden der ersten Iteration und die Maxima der Verteilung
    \item Wer den kurativen Reservierungspreis senken will, muss deshalb an der Regelleistung ansetzen
\end{itemize}

\noindent\textbf{Absatz 3, die Wirksamkeit gegenüber dem Engpassmanagement.}
\begin{itemize}
    \item Übergang, ob sich die Vorhaltung lohnt, bemisst sich am bestehenden Engpassmanagement
    \item Das Tagesmittel des kurativen Reservierungspreises bleibt fast das ganze Jahr unter dem Arbeitspreis des präventiven Redispatch \cite{bundesnetzagentur_smard_engpassmanagement_2026}
    \item Beide Größen tragen dieselbe Dimension und meinen Verschiedenes, denn die eine vergütet Vorhaltung und die andere bewegte Arbeit
    \item Im Winter trifft der größte Engpassmanagementbedarf auf die günstigste Reservierung, was der stärkste Befund für die Beschaffung ist
    \item Die Mittagsstunden bilden die Ausnahme, denn dort fällt der größte Bedarf an Reduzierung der Einspeisung mit der teuersten Ladereservierung zusammen
    \item Ob eine kurative Maßnahme eine Redispatchmaßnahme ersetzt, hängt an der Zahl der vorgehaltenen Stunden je Abruf, die diese Arbeit nicht bestimmt
\end{itemize}

\noindent\textbf{Absatz 4, der Zuschnitt des Produkts und der Vergütungsrahmen.}
\begin{itemize}
    \item Übergang, über die Kosten des \ac{ÜNB} entscheidet neben dem Preis der Zuschnitt des Produkts
    \item Ein Produkt mit einer Zeitscheibe von einem Tag bindet die Stunden unvollständig und kostet dennoch beinahe den vollen Preis
    \item Die kurative Reservierung ist deshalb je Stunde auszuschreiben
    \item Eine Abrufdauer von einer halben Stunde ist fast zum Preis einer Viertelstunde zu haben und verdoppelt die Zeit für eine Ablösung der Maßnahme
    \item Die Vergütungsstruktur der geltenden Festlegung erfasst den ermittelten Preis nicht, denn sie wählt zwischen Werteverbrauch und Opportunität, statt beide zu addieren \cite{bundesnetzagentur_festlegung_2024}
    \item Eine Vergütung der kurativen Vorhaltung verlangt deshalb eine eigene Struktur, zumal für die kurative Reservierung bislang kein Vergütungsrahmen besteht
\end{itemize}

% ---------------------------------------------------------------
% TITEL UMGESTELLT am 24.09.2026: ein acs am Anfang einer Ueberschrift
% zerlegt den Eintrag im Inhaltsverzeichnis. Der Titel beginnt deshalb
% mit einem Wort. Zurueckgenommene Fassung: BESS im Engpassmanagement
% und am Markt
\section{Die Rolle der \acs{BESS} im Engpassmanagement und am Markt}
\label{sec:bess_markets}
% Ziel rund 2,5 Seiten, vier Absaetze. Der Abschnitt, den der Verfasser
% am 23.09.2026 ausdruecklich verlangt hat, um den Zusammenhang von
% Batteriespeichern und den kurzfristigen Befunden erweitert.

\noindent\textbf{Absatz 1, warum ein \ac{BESS} für die kurative Vorhaltung geeignet ist.}
\begin{itemize}
    \item Übergang, unter den Technologien ist das \ac{BESS} für die kurative Vorhaltung besonders geeignet
    \item Es erreicht die Reaktionszeitklasse von etwa zwei Minuten ohne Anfahrvorgang \cite{innosys_2030_gesamtverbund_innovationen_2021}
    \item Es stellt beide Richtungen aus derselben Anlage, sodass eine bilanziell ausgeglichene Vorhaltung an einem Netzknoten möglich ist
    \item Sein Energieinhalt je Leistung genügt für die im Produkt vorgesehene Abrufdauer
    \item Seine Präqualifikation kann auf der Regelleistung aufsetzen, sodass der zusätzliche Aufwand gering bleibt
    \item Die kurative Vorhaltung trifft damit auf eine Technologie, deren Bestand ohnehin zunimmt
\end{itemize}

\noindent\textbf{Absatz 2, der Preis eines \ac{BESS} hängt an den kurzfristigen Märkten.}
\begin{itemize}
    \item Übergang, der kurative Reservierungspreis eines \ac{BESS} folgt aus seinen kurzfristigen Erlösquellen
    \item Die Vorhaltung von \ac{aFRR}-Leistung trägt den größten Teil seines Erlöses, der kontinuierliche Intraday-Handel den zweitgrößten
    \item Den Leistungspreis der \ac{aFRR} kennt ein Betreiber bei der Gebotsabgabe, die Handelsspannen des folgenden Tages kennt er nicht
    \item Gerade die Handelsspannen tragen die teuren Stunden, sodass ein kuratives Gebot unter Unsicherheit entsteht
    \item Wer zu niedrig bietet, bindet seine Leistung unter Wert, und wer zu hoch bietet, erhält keinen Zuschlag
    \item Die kurative Reservierung gibt dem Betreiber zugleich eine Möglichkeit, seinen Erlös allein durch eine klügere Verteilung auf die Märkte zu steigern
\end{itemize}

\noindent\textbf{Absatz 3, wo die Grenze der Eignung liegt.}
\begin{itemize}
    \item Übergang, die Eignung hat drei Grenzen
    \item Die erste ist der Energieinhalt, denn eine längere Abrufdauer verlangt ein breiteres Ladezustandsband und verteuert die Vorhaltung
    \item Die zweite ist die Doppelbelegung dieses Bandes, das die Vorhaltung von Regelleistung bereits beansprucht
    \item Die dritte ist die örtliche Wirksamkeit, die am Netzknoten hängt und nicht an der Technologie
    \item Ein \ac{BESS} am falschen Netzknoten ist für die Engpassbehebung wertlos, so schnell es auch reagiert
    \item Über den Nutzen entscheidet damit die Standortwahl ebenso wie die technische Eignung
\end{itemize}

\noindent\textbf{Absatz 4, die künftige Rolle im Strommarkt und bei den Systemdienstleistungen.}
\begin{itemize}
    \item Übergang, mit dem Bestand verschiebt sich, wofür ein \ac{BESS} eingesetzt wird
    \item Der Bedarf an Regelleistung folgt der Dimensionierung des Systems und wächst nicht mit dem Bestand an \ac{BESS}
    \item Für die \ac{FCR} ist dieser Punkt erreicht, denn die präqualifizierte Batterieleistung übersteigt den deutschen Bedarf \cite{celi_cortes_m5use_2024, bundesnetzagentur_monitoringbericht_2026}
    \item Die Vorhaltung von \ac{aFRR}-Leistung ist der nächste Markt, der diesen Punkt erreicht
    \item Sinkt ihr Preis, so sinkt der kurative Reservierungspreis mit ihm, ohne dass sich sein Verhältnis zur Marktopportunität ändert
    \item Die Beschlusskammer hat die Behandlung von \ac{BESS} einem späteren Hinweis vorbehalten, sodass die Regelungslücke dokumentiert ist \cite{bundesnetzagentur_festlegung_2024}
\end{itemize}

% ---------------------------------------------------------------
\section{Umsetzung, systemweite Ausrollung und Forschungsbedarf}
\label{sec:implementation_outlook}
% Ziel rund 3 Seiten, vier Absaetze. Der Abschluss, den der Verfasser am
% 23.09.2026 verlangt hat: Aufwand der Umsetzung, systemweite Ausrollung,
% Abwaegung und daraus der Forschungsbedarf.

\noindent\textbf{Absatz 1, was der geltende Rahmen zulässt.}
\begin{itemize}
    \item Übergang, vor der Umsetzung steht der geltende Rahmen
    \item §\,13a Abs.\,2 \ac{EnWG} entschädigt den durchgeführten Eingriff und schreibt eine kostenbasierte Abrechnung vor \cite{bundesministerium_der_justiz_energiewirtschaftsgesetz_2025}
    \item Eine marktbasierte Beschaffung des Redispatch ist damit im geltenden Recht ausgeschlossen
    \item Redispatch~2.0 ist ein Duldungs- und Entschädigungsregime, während die kurative Reservierung ein Gebot des Betreibers voraussetzt \cite{bundesnetzagentur_festlegung_2024}
    \item Der entworfene kurative Marktmechanismus setzt deshalb eine Änderung des §\,13a \ac{EnWG} voraus
    % NEU EINZUFUEHRENDER BEGRIFF, Beleg noch zu pruefen.
    \item Der kapazitätsbasierte Redispatch beschafft im Kern dieselbe Größe und unterscheidet sich nur in Auslöser und Reaktionszeit, sodass ein Anschluss an eine bestehende Beschaffungsform möglich ist
\end{itemize}

\noindent\textbf{Absatz 2, der Aufwand der Umsetzung.}
\begin{itemize}
    \item Übergang, neben dem Rahmen steht der Aufwand beim \ac{ÜNB} und beim Betreiber
    \item Der \ac{ÜNB} muss die bezuschlagten Gebote zu einem kurativen Maßnahmenset zusammenfassen und in Betriebsplanung, Scharfschaltung und Auslösung führen
    \item Die Vergütung der Vorhaltung nimmt dem Betreiber den Anreiz, das Volumen eines Eingriffs zu vergrößern
    \item An dessen Stelle tritt der Anreiz, Leistung anzubieten, die nicht wirklich frei ist, sodass aus einer Mengenfrage eine Nachweisfrage wird
    \item Ein Verfahren zum Nachweis der Verfügbarkeit ist deshalb konstitutiver Bestandteil des Mechanismus und keine Frage der Ausführung
    \item Die Präqualifikation kann auf der Regelleistung aufsetzen, sodass sich der zusätzliche Aufwand auf die kurative Zusage beschränkt
\end{itemize}

\noindent\textbf{Absatz 3, die systemweite Ausrollung.}
\begin{itemize}
    \item Übergang, ob sich der Mechanismus auf das ganze System übertragen lässt, ist eine eigene Frage
    \item Diese Arbeit rechnet eine einzelne Anlage und trägt deshalb keine systemweite Mengenaussage
    \item Eine Hochrechnung setzte voraus, dass die Wirksamkeit eines Akteurs eine feste Eigenschaft seines Netzknotens ist
    \item Die Wirksamkeit ist keine feste Eigenschaft, denn die vorherrschende Erzeugungssituation verschiebt sie über die Jahre \cite{bundesnetzagentur_smard_engpassmanagement_2026}
    \item Die kurative Systemführung ergänzt das präventive Engpassmanagement und ersetzt es nicht
    \item Ihr Nutzen liegt darin, dem \ac{ÜNB} neben dem präventiven Eingriff ein zweites Werkzeug zu geben
\end{itemize}

\noindent\textbf{Absatz 4, was zu untersuchen bleibt.}
\begin{itemize}
    \item Übergang, aus den Grenzen dieser Arbeit folgt der weitere Forschungsbedarf
    \item Die Wahrscheinlichkeit eines Abrufs und damit der Optionswert der Bindung sind zu bestimmen
    \item Eine Hochrechnung auf die Systemebene verlangt ein Modell, das das Netz und die Wirksamkeit je Netzknoten führt
    \item Das Bietverhalten und der Wettbewerb mehrerer Anbieter an einem Netzknoten sind abzubilden
    \item Aus Projekten im Echtbetrieb ist zu lernen, wie häufig und wie lange eine kurative Maßnahme abgerufen wird und wie sich die Verfügbarkeit im Betrieb nachweisen lässt \cite{innosys_2030_gesamtverbund_innovationen_2021, tennet_tso_gmbh_pilotbetrieb_2025}
    \item Erst damit bekommt das Engpassmanagement neben dem präventiven Eingriff ein zweites Werkzeug, mit dem sich das Übertragungsnetz sicher betreiben lässt
\end{itemize}

% ===============================================================
% ARCHIV, nur zum Nachschlagen. Nicht wieder aktivieren.
%
% ERSTENS die am 24.09.2026 vormittags ausformulierte
% Kapiteleinleitung der Gliederung mit acht Abschnitten. Sie ist
% Fliesstext gewesen und bleibt deshalb nach CLAUDE.md Abschnitt 10
% erhalten. Ersetzt worden ist sie am selben Tag nachmittags, weil der
% Verfasser die Einleitung ohne Zahlen und auf den neuen Aufbau
% bezogen verlangt hat.
% ===============================================================
%Die Bewertung der Ergebnisse setzt eine Bestandsaufnahme voraus.
%In Kapitel~\ref{ch:results} wird für jede Stunde des Jahres 2025 je Richtung ein kurativer Reservierungspreis bestimmt.
%Ein steigender vorgegebener Preis löst dabei nacheinander die Vermarktung an den übrigen Märkten ab.
%Welcher Markt zuerst weicht, folgt der Preisstruktur des jeweiligen Tages.
%Über das Jahr bindet die \ac{aFRR}-Leistung an jedem der fünf untersuchten Tage den größten Teil der Leistung.
%Auf die \ac{aFRR} entfallen 76,7 Prozent des Bruttoerlöses, sodass die kurative Reservierung vor allem die Vorhaltung von Regelleistung verdrängt.
%Die Höhe dieses Preises unterscheidet sich von Stunde zu Stunde.
%Der Median beträgt 12,22\,\EurMWh{} in der Entlade- und 9,85\,\EurMWh{} in der Laderichtung.
%Die Verteilung ist in beiden Richtungen rechtsschief, denn das arithmetische Mittel liegt beim Zwei- bis Zweieinhalbfachen des Medians.
%Die Zahlung des \ac{ÜNB} bei vollständiger Bindung beträgt 420,0\,\TsdEurMWa{} gegen einen Referenzerlös von 340,6\,\TsdEurMWa{}.
%Die Preise der ersten Iteration binden die Stunden unvollständig und kosten dabei bereits 97,6 Prozent der vollen Zahlung.
%Für die Beschaffung zählt daneben, wann der kurative Reservierungspreis niedrig liegt.
%Über den Tag schwankt der Median der Summe beider Richtungen um den Faktor 8,3 und über die Monate um 4,2.
%Die Ladereservierung ist mittags teuer und bildet von März bis Oktober ein zusammenhängendes Band, während die Entladereservierung über den Vormittag und den Abend streut.
%Im Winter kostet die Reservierung rund 40 Prozent weniger als im Sommer, und dort liegt zugleich der größte Engpassmanagementbedarf.
%Das Tagesmittel bleibt an mindestens 362 von 365 Tagen unter dem Arbeitspreis des präventiven Redispatch von 101 Euro je bewegter Megawattstunde.
%Woran dieses Ergebnis hängt, beantworten vier Sensitivitäten.
%Die Annahmen über die \ac{aFRR} tragen es am stärksten, denn ohne die \ac{aFRR} fällt der Median auf 28 Prozent in der Entlade- und auf 43 Prozent in der Laderichtung.
%Unter freier Lieferung der \ac{aFRR}-Arbeit steigt der Median um 54 und 110 Prozent, was für eine einzelne Anlage von 100\,\si{MW} nicht zu erwarten ist.
%Eine verdoppelte Handelsspanne am \ac{IDC} hebt den Median um gut ein Drittel und die Maxima um das Doppelte.
%Eine Viertelung der vorgehaltenen Energie je Abruf senkt den Median dagegen nur um 4,0 und 12,7 Prozent.
%Die Maxima reagieren auf jede Annahme stärker als der Median, sodass die Spitze der Verteilung weniger belastbar ist als ihr Hauptfeld.
%Die vierte Sensitivität prüft statt des Optimierungsmodells eine Annahme über die künftige Marktlage.
%Sinken die Preise der \ac{aFRR} und die Handelsspanne am \ac{IDC} auf 80 Prozent, so fällt der Markterlös um 17,0 Prozent.
%Die Zahlung bei vollständiger Verdrängung fällt um 16,0 Prozent.
%Der kurative Reservierungspreis gibt damit weniger nach als der Markterlös, aus dem er folgt.
%Das Niveau sinkt mit dem Markt, während das Verhältnis beider bleibt.
%Diese Ergebnisse sind einzuordnen, zu bewerten und auf ihre Folgen zu prüfen.
%Zu Anfang steht deshalb, was der ermittelte Preis aussagt und wie er sich in den geltenden Rahmen fügt.
%Darauf folgt die Bewertung des kurativen Marktmechanismus am Anforderungskatalog aus Abschnitt~\ref{sec:requirements_catalogue}.
%Anschließend werden die Folgen für die Betriebsführung und für den Zuschnitt des Produkts gezogen und die Belastbarkeit des Ergebnisses geprüft.
%Zuletzt stehen die Entwicklungen an den Märkten, die Zukunft der \acs{BESS} im Engpassmanagement und die kurative Systemführung im Allgemeinen, aus der der weitere Forschungsbedarf folgt.
%% ---------------------------------------------------------------

% Die Stichpunkte der Gliederung mit acht Abschnitten stehen im Git
% unter dem Commit e8c4012 und sind hier nicht wiederholt, weil sie
% kein Fliesstext waren.
% ===============================================================
% ALTES STICHPUNKTGERUEST vom 18.08.2026, am 24.09.2026 durch die
% Fassung oben ersetzt. Vollstaendig auskommentiert erhalten, weil es
% den Belegstand jedes Punktes traegt, naemlich [belegt], [eigen] und
% [daten] samt Fundstelle. Die Kennungen B1 bis G4 sind hier
% nachzulesen. NICHT WIEDER AKTIVIEREN, die Gliederung ist am
% 24.09.2026 vom Verfasser durch Fassung A ersetzt.
% ===============================================================
%%%%%%%%%%%%%%%%%%%%%%%%%%%%%%%%%%%%%%%%%%%%%%%%%%%%%%%%%%%%%%%%%%%%%%%
%%%                       IAEW LaTeX Template                        %%
%%%%%%%%%%%%%%%%%%%%%%%%%%%%%%%%%%%%%%%%%%%%%%%%%%%%%%%%%%%%%%%%%%%%%%%
%
% \chapter{Bewertung und Diskussion}
% \label{ch:discussion}
%% ===============================================================
%% Zielumfang Kapitel 5: 12 Seiten
%%   5.1 Einordnung in den regulatorischen Rahmen ............. 3
%%   5.2 Bewertung anhand des Anforderungskatalogs ............ 3
%%   5.3 Uebertragbarkeit auf Systemebene ..................... 2
%%   5.4 Kritische Wuerdigung der Modellannahmen .............. 3
%%   5.5 Handlungsempfehlungen ................................ 1
%%
%% STAND DIESER FASSUNG: Die Punktelisten sind aus der Sammlung
%% abgrenzung_diskussion_sammlung.md uebertragen. Die Kennungen B1, C1,
%% D3 und so weiter verweisen dorthin, damit Herkunft und Belegstand
%% jedes Punktes nachpruefbar bleiben. Die Kennungen sind vor Abgabe zu
%% entfernen.
%%
%% BELEGSTAND, in den Kommentaren jeweils mitgefuehrt:
%%   [belegt]  Fundstelle vorhanden, direkt verwendbar
%%   [eigen]   in der Bearbeitung entstandene Argumentation, vor der
%%             Aufnahme selbst zu pruefen und in eigenen Worten zu
%%             schreiben
%%   [daten]   eigene Datenauswertung, Zahl aus dem Repository belegen
%%
%% STILVORGABE wie in Kapitel 1 und 2: vollstaendige Saetze, keine
%% Doppelpunkte, keine Gedankenstriche, keine Semikola im Fliesstext.
%% Die Vorfassung dieser Datei enthielt in zwei Listenpunkten einen
%% Gedankenstrich und ein Semikolon. Beide sind aufgeloest.
%%
%% BEGRIFFLICHKEIT FESTGELEGT, Stand 18.08.2026: Die gesuchte Groesse
%% heisst kurativer Reservierungspreis. Aufgegeben sind
%% "Mindestverguetung", "Verguetungsbereich" und "Untergrenze der
%% Verguetung". Ein Reservierungspreis ist definitionsgemaess der Betrag,
%% unter dem ein Anbieter nicht verkauft, und traegt die Aussage damit
%% ohne Hilfskonstruktion. Grenzpreis und nicht Schwellenpreis, weil die
%% Anlage 2 zur Festlegung diesen Begriff fuehrt.
%%
%% WEITERE ENTSCHEIDUNGEN VOM 18.08.2026:
%%   1  Der modifizierte Weber-Ansatz wird nicht nachgerechnet. Punkt F6
%%      traegt deshalb nur noch die methodische Begruendung und keine
%%      Zahlen. Die Begruendung selbst steht in Abschnitt 3.4.1.
%%   2  Die Dualvariablen sind aus der Methodik entfallen. Der frueher
%%      hier stehende Punkt F9 zur Gewinnung der Grenzpreise aus dem
%%      Problem mit fixierten Zusagen ist damit gegenstandslos und
%%      entfernt.
%%   3  Zwei Pruefnotizen sind erledigt und entfernt, naemlich der
%%      Hinweis, Kapitel 1 und 2 beschrieben die Verguetungsstruktur noch
%%      additiv, und der Hinweis, Redispatch 2.0 sei noch als
%%      marktbasiert charakterisiert. Beides trifft auf die vorliegende
%%      Fassung nicht mehr zu.
%%
%% ABGRENZUNG GEHOERT NICHT IN DIESES KAPITEL: Die Punkte A1 bis A5 der
%% Sammlung sagen, was die Arbeit nicht untersucht. Das sind Festlegungen
%% und keine Verzerrungen des Modells. Sie gehoeren in die Systemgrenzen
%% in Abschnitt 3.3.1 und werden hier nur in ihrer Konsequenz aufgegriffen,
%% naemlich A2 und A3 in Abschnitt 5.3 und A5 in Abschnitt 5.2.
%% Der Kommentarblock am Ende von chapter_1.tex ist entsprechend nach
%% chapter_3.tex zu uebertragen und nicht nach chapter_6.tex.
%% ===============================================================
%
%% ---------------------------------------------------------------
% \section{Einordnung in den regulatorischen Rahmen}
% \label{sec:regulatory_discussion}
%% [Ziel: 3 Seiten]
%%
%% ARGUMENTATIONSBOGEN DES ABSCHNITTS: Der ermittelte Verguetungsbereich
%% wird an eine Systematik gehalten, die Vorhaltung nicht kennt. Daraus
%% folgt, dass er innerhalb des geltenden Rahmens nicht abbildbar ist.
%% Das ist die Einordnung, nicht die Bewertung. Die Bewertung gegen die
%% eigenen Anforderungen folgt in Abschnitt 5.2.
%%
%% PRUEFEN: Die Verweise \ref{sec:market_development} und
%% \ref{sec:requirements_catalogue} setzen die Labels aus Kapitel 2
%% voraus. Falls die Abschnitte 3.1.1 und 3.1.2 nach Kapitel 3 wandern,
%% sind beide Verweise nachzuziehen.
%
% \begin{itemize}
%     % B1 [belegt] Kern der regulatorischen Einordnung.
%     \item Verhältnis des ermittelten kurativen Reservierungspreises zur Systematik des §\,13a \ac{EnWG}, die den durchgeführten Eingriff entschädigt und keine Vorhaltung vergütet \cite{bundesministerium_der_justiz_energiewirtschaftsgesetz_2025, bundesnetzagentur_festlegung_2024}
%     % B2 [belegt fuer die Struktur, eigen fuer die Folgerung auf BESS]
%     % TRAGENDER BEFUND. Die Festlegung bildet Erzeugungsauslagen plus
%     % max(Werteverbrauch, Opportunitaet). Der Weber-Optionswert kann
%     % damit in die Zahlung ueberhaupt nicht eingehen, wenn der
%     % Werteverbrauch groesser ist. Fuer Batteriespeicher ist genau das
%     % zu erwarten.
%     % ERLEDIGT am 18.08.2026: Abschnitt 2.3.2 fuehrt die Struktur korrekt
%     % als Auswahlverhaeltnis, Kapitel 1 beschreibt sie an keiner Stelle.
%     \item Nicht additive Struktur der Vergütung nach geltender Festlegung und die Folge, dass der Optionswert für Batteriespeicher voraussichtlich nicht zahlungswirksam wird \cite{bundesnetzagentur_festlegung_2024}
%     % B3 [belegt] ERLEDIGT am 18.08.2026: Abschnitt 2.3.2 fuehrt
%     % Redispatch 2.0 als Duldungs- und Entschaedigungsregime.
%     \item Abgrenzung des entworfenen Mechanismus gegenüber Redispatch~2.0 als Duldungs- und Entschädigungsregime \cite{bundesnetzagentur_festlegung_2024}
%     % B4 [belegt] Die Beschlusskammer 8 hat die Behandlung von
%     % Batteriespeichern einem spaeteren Hinweis vorbehalten. Damit ist
%     % die Forschungsluecke ein dokumentierter offener Punkt.
%     % Herleitung steht in Abschnitt 3.1.3, hier nur die Einordnung des
%     % eigenen Ergebnisses in diesen offenen Punkt.
%     \item Einordnung des Ergebnisses in den von der Beschlusskammer ausdrücklich offengehaltenen Regelungsbedarf für Speicheranlagen \cite{bundesnetzagentur_festlegung_2024}
%     % B5 [daten] Regimebruch 01.10.2025. Der innere Wert des
%     % Weber-Ansatzes hat mit der Viertelstundenauktion seine
%     % Bemessungsgrundlage verloren.
%     % PRUEFEN: Kennzahl aus dem Repository belegen, RMS(ID1 minus IDA1)
%     % von 46,8 auf 28,4 Euro je MWh.
%     \item Folgen der Verlagerung von Handel aus der Intraday-Auktion in den Day-Ahead-Markt für die Bemessungsgrundlage des modifizierten Weber-Ansatzes \cite{ffe_saegezahn_2026}
%     % B6 [eigen, gestuetzt auf mahgoub_wie_2025] Argument fuer
%     % Vorhalteverguetung. Eine arbeitsbezogene Verguetung waere an eine
%     % Bemessungsgrundlage gebunden, die sich unter ihr veraendert.
%     \item Rückgriff auf Abschnitt \ref{sec:market_development}, wonach ein kurativer Markt in eine Marktentwicklung eingreift, die selbst in Bewegung ist, und was daraus für die Vergütung von Vorhaltung anstelle von Arbeit folgt \cite{mahgoub_wie_2025}
%     % B7 [belegt]
%     \item Vergleichbarkeit mit der Behandlung von Anlagen zur Erzeugung aus erneuerbaren Energien \cite{iaew_abschlussbericht_2023, iaew_folgestudie_2024}
% \end{itemize}
%
%% ---------------------------------------------------------------
% \section{Bewertung anhand des Anforderungskatalogs}
% \label{sec:requirements_assessment}
%% VORMERKUNG vom 11.09.2026 aus 3.1.2, Stichpunkt zur Einfuegung in die
%% Betriebsplanung. Bei A8 ist hier der Vorbehalt zu nennen, dass sich der
%% Aufwand der Einfuegung in die Betriebsprozesse von aussen nicht beurteilen
%% laesst, weil er die internen Ablaeufe der UENB betrifft. Der Verfasser will
%% A8 deshalb allgemeiner bewerten. Die Nummern der Anforderungen in diesem
%% Kapitel folgen teils noch dem Katalog mit sechs Anforderungen.
%% [Ziel: 3 Seiten]
%%
%% Systematischer Abgleich gegen Abschnitt \ref{sec:requirements_catalogue}.
%% Zusaetzlich zu den dort formulierten Anforderungen sind drei Punkte zu
%% bewerten, die sich erst aus der Modellierung ergeben.
%%
%% GAMING: Dieser Abschnitt traegt die ausfuehrliche Behandlung. Kapitel 1
%% zeigt die Auseinandersetzung nur an, Abschnitt 2.3.2 stellt sie im
%% Verguetungsregime dar. Die Punkte D1 und D2 der Sammlung gehoeren nach
%% 2.3.2, die Punkte D3 bis D5 hierher.
%% VORAUSSETZUNG: Der BMWi-Abschlussbericht von Hirth, Schlecht, Maurer
%% und Tersteegen 2019 muss in der literature.bib stehen, sonst ist nur
%% die ablehnende Position der Beschlusskammer belegt und die Diskussion
%% steht einseitig. PDF unter neon.energy. Bis dahin sind die
%% Platzhalterschluessel unten nicht aufloesbar.
%
% \begin{itemize}
%     % C1 [belegt] STARKER PUNKT. Die FCR-Praequalifikationsbedingungen
%     % definieren bereits einen reservierten SoC-Arbeitsbereich, die
%     % Schnabelkurve, samt Fuenfzehnminutenkriterium. Die kurative
%     % Vorhaltung ist damit der zweite Anspruch auf dasselbe Band.
%     \item Konkurrenz der kurativen Vorhaltung zu dem bereits durch die Präqualifikationsbedingungen reservierten Arbeitsbereich des Ladezustands \cite{ubertragungsnetzbetreiber_deutschland_praqualifikationsverfahren_2024}
%     % C2 [belegt] BK6-17-234 vom 02.05.2019. Praezedenzfall dafuer, dass
%     % die Bandbreite eine regulatorische Entscheidung mit unmittelbarer
%     % Erloeswirkung ist. Traegt die Argumentation zur Vorhaltedauer.
%     \item Bandbreite und Vorhaltedauer als regulatorische Entscheidung mit unmittelbarer Erlöswirkung
%     % D3 [eigen] ARGUMENT FUER DEN EIGENEN ENTWURF, steht bislang
%     % nirgends in der Arbeit. Haengt die Zahlung an der
%     % bereitgehaltenen Leistung und nicht am abgerufenen Volumen, so
%     % entfaellt der Anreiz, das Volumen zu vergroessern. Der Akteur
%     % verdient an der Verfuegbarkeit und nicht am Eingriff.
%     % Vor der Aufnahme gegen die Literatur zu Reliability Options
%     % pruefen, damit das Argument nicht als eigene Entdeckung steht,
%     % falls es dort bereits formuliert ist.
%     \item Wirkung der Vergütung von Vorhaltung anstelle von Arbeit auf den Anreiz zu engpassverstärkendem Bietverhalten
%     % D4 [eigen, gestuetzt auf C1] STAERKSTER EIGENER BEITRAG in diesem
%     % Block, weil er das Thema aus der Abwehrhaltung holt. Das
%     % Anreizproblem verschwindet nicht, es wandert. Eine
%     % Vorhalteverguetung schafft den Anreiz, Kapazitaet anzubieten, die
%     % nicht wirklich frei ist. Aus der Volumenmanipulation wird eine
%     % Nachweisfrage, und die Pruefbarkeit der Verfuegbarkeit wird zur
%     % konstitutiven Anforderung an den Mechanismus.
%     \item Verlagerung des Anreizproblems von der Menge des Eingriffs auf den Nachweis der Verfügbarkeit und die daraus folgende Anforderung an die Prüfbarkeit
%     % D5 [eigen] Die Opportunitaetskosten sind private Information des
%     % Bieters. Bei vielen geeigneten Anbietern disziplinieren sich
%     % Ueberzeichnungen gegenseitig, an einem Knoten mit wenigen
%     % wirksamen Anlagen nicht. Die Lage des Preises im
%     % Reservierungspreis ist damit keine Modellgroesse, sondern eine
%     % Frage der Anbieterzahl. Der Satz dazu steht bereits in
%     % Abschnitt 1.2 und ist hier aufzuloesen.
%     \item Lokale Marktmacht bei geringer Zahl wirksamer Anbieter an einem Netzknoten und die Folge für die Lage des Preises oberhalb des ermittelten kurativen Reservierungspreises
%     % A5 [folgt aus A2] Offen sagen, dass die Arbeit das Bietverhalten
%     % nicht modelliert und die Aussagen dieses Abschnitts daher
%     % qualitativ bleiben.
%     \item Grenzen dieser Bewertung, da das Modell weder eine Bietfunktion noch einen Wettbewerb mehrerer Anbieter abbildet
% \end{itemize}
%
%% ---------------------------------------------------------------
% \section{Übertragbarkeit auf Systemebene}
% \label{sec:transferability}
%% [Ziel: 2 Seiten]
%%
%% Hier wird das Skalenproblem offen adressiert, statt darauf zu warten,
%% dass es im Kolloquium kommt. Eine Einzelanlagenoptimierung traegt
%% keine systemweite Volumenaussage. Voraussetzungen und Grenzen einer
%% Hochrechnung benennen.
%
% \begin{itemize}
%     % E1 [eigen, folgt aus A2] Konsequenz der Einzelanlagenperspektive.
%     % Der Widerspruch zwischen Forschungsfrage und Methodik ist in
%     % Kapitel 1 bereits aufgeloest, hier steht die Folge fuer die
%     % Reichweite der Ergebnisse.
%     \item Voraussetzungen, unter denen eine Hochrechnung des Einzelanlagenergebnisses auf die Systemebene zulässig wäre, und Gründe, warum sie hier nicht vorgenommen wird
%     % E2 [belegt] Die Wirksamkeit eines Akteurs ist keine feste
%     % Eigenschaft seines Standorts. Die Verschiebung von der
%     % Nord-Sued- auf die Ost-West-Achse zeigt das, PV-Redispatch plus
%     % 94 Prozent, Offshore minus 27 Prozent im Jahr 2025.
%     \item Abhängigkeit der Wirksamkeit eines Akteurs von der jeweils vorherrschenden Erzeugungssituation und die Folge für eine Hochrechnung über Netzknoten \cite{bundesnetzagentur_smard_engpassmanagement_2026}
%     % A3 [Festlegung] Konsequenz der unterstellten netztechnischen
%     % Wirksamkeit fuer die Reichweite der Aussage.
%     \item Reichweite der Ergebnisse unter der Annahme gegebener netztechnischer Wirksamkeit der Maßnahme
% \end{itemize}
%
%% ---------------------------------------------------------------
% \section{Kritische Würdigung der Modellannahmen}
% \label{sec:model_critique}
%% [Ziel: 3 Seiten]
%%
%% BEIDE Verzerrungsrichtungen darstellen, sonst wirkt der Abschnitt
%% einseitig.
%%   - Vollstaendige Preiskenntnis UEBERSCHAETZT die Erloesdifferenz.
%%     Weber raeumt im Fazit seines Gutachtens dasselbe fuer sein
%%     Verfahren ein, naemlich dass die Vernachlaessigung
%%     zeituebergreifender Restriktionen tendenziell zur Ueberschaetzung
%%     fuehrt.
%%   - Vollstaendige Preiskenntnis UNTERSCHLAEGT den Optionswert
%%     vollstaendig, bei sigma gleich null ist er identisch null, weshalb
%%     er separat hinzugerechnet werden muss.
%% Ergaenzend: Der Unterschied zwischen Erloespotenzial und tatsaechlich
%% realisiertem Erloes ist unabhaengig belegt \cite{mahgoub_wie_2025}.
%%
%% ENTFALLEN gegenueber der Vorfassung: Der Punkt "Betrachtung eines
%% einzelnen Tages" ist ueberholt. Mit dem Regimebruch vom 01.10.2025 ist
%% der Auswertungszeitraum das vierte Quartal 2025, siehe F5.
%% UEBERHOLT am 11.09.2026, Entscheidung des Verfassers. Die Arbeit rechnet das
%% ganze Jahr 2025.
%
% \begin{itemize}
%     % F1 und F2 [methodisch zwingend, Weber belegt] Die beiden
%     % Verzerrungen laufen gegeneinander. Das ist der Grund fuer die
%     % hybride Ex-post-Rechnung und als solcher zu benennen. Der
%     % ermittelte Wert ist damit eine konservative Obergrenze, was keine
%     % Schwaeche, sondern eine saubere Eigenschaft des Ergebnisses ist.
%     \item Vollständige Preiskenntnis und ihre beiden gegenläufigen Verzerrungsrichtungen sowie die daraus folgende Einordnung des Ergebnisses als konservative Obergrenze \cite{weber_gutachten_2015, mahgoub_wie_2025}
%     % F3 [folgt aus A2 und A3]
%     \item Einzelanlagenbetrachtung ohne Abbildung des Netzes
%     % F4 [eigen] Ob der Fehlerfall eintritt, ist eine
%     % Wahrscheinlichkeit und keine Fahrplangroesse. Das Problem hat
%     % Optionscharakter. Damit ist der Weber-Ansatz nicht nur
%     % Kontrastmodell, sondern der methodisch naheliegende naechste
%     % Schritt, was den Ausblick in Abschnitt 6.2 traegt.
%     \item Deterministischer Abruf anstelle einer Abrufwahrscheinlichkeit
%     % F5 [daten] NEU. Der Regimebruch vom 01.10.2025 macht das vierte
%     % Quartal 2025 zum einzigen gueltigen Auswertungsfenster fuer die
%     % gewaehlte Architektur. Das ist eine kurze Zeitreihe, die
%     % Saisonalitaet ist damit nicht abgedeckt.
%     % UEBERHOLT am 11.09.2026, Entscheidung des Verfassers. Die Arbeit rechnet
%     % das ganze Jahr 2025, die Saisonalitaet ist damit abgedeckt. Alte Fassung:
%     %    \item Beschränkung des Auswertungszeitraums auf das vierte Quartal 2025 und die daraus folgende fehlende Abdeckung saisonaler Effekte
%     % VORSCHLAG, vom Verfasser zu bestaetigen. An die Stelle koennte der
%     % Regimebruch innerhalb des Zeitraums treten, naemlich drei Quartale mit
%     % Stundenprodukt am Day-Ahead und eines mit Viertelstundenprodukt.
%     % F6 [eigen] GEAENDERT am 18.08.2026. Die Zahlen sind entfernt, weil
%     % sie eine Nachrechnung voraussetzen, die nicht stattfindet. Die
%     % drei methodischen Gruende stehen in Abschnitt 3.4.1, hier steht ihre
%     % Folge fuer die Reichweite des Ergebnisses, naemlich dass der
%     % ermittelte Wert rein opportunitaetskostenbasiert ist und keinen
%     % Optionswert enthaelt. Abschnitt 2.3.3 nennt daneben sechs
%     % inhaltliche Gruende gegen die Anwendbarkeit auf Batteriespeicher.
%     % Beide Ebenen nicht vermengen.
%     \item Folge der in Abschnitt \ref{sec:weber_exclusion} begründeten Abgrenzung, nämlich dass der ermittelte kurative Reservierungspreis rein opportunitätskostenbasiert ist und keinen Optionswert enthält \cite{weber_gutachten_2015, bundesnetzagentur_festlegung_2024}
%     % ENTFALLEN am 18.08.2026: Der Punkt F9 zur Gewinnung der Grenzpreise
%     % aus dem Problem mit fixierten Binaervariablen ist gegenstandslos,
%     % weil die Dualvariablen aus der Methodik entfallen sind. Sie wurden
%     % allein fuer den Vergleich mit dem Weber-Grenzpreisband gebraucht,
%     % und dieser Vergleich findet nicht statt. Nicht wieder aufnehmen,
%     % ohne zugleich Abschnitt 3.3.3 umzustellen.
%     % F10 [methodisch] NEU. Die Verifikation ueber den Revenue Index der
%     % ISEA Battery Charts vergleicht ein Modell mit vollstaendiger
%     % Preiskenntnis gegen ein Rolling-Intrinsic-Verfahren und einen
%     % engeren gegen einen weiteren Marktumfang. Beide Unterschiede sind
%     % hier als Grenze der Aussage zu benennen.
%     \item Reichweite der Verifikation über einen extern berechneten Erlösindex mit abweichender Vorausschau und abweichendem Marktumfang
%     % F11 [eigen] NEU am 12.09.2026 nach der Anweisung des Verfassers. Der
%     % Vektor der 2. Iteration ist ein koordinatenweises Minimum, kein globales.
%     % Keine einzelne Stunde laesst sich um die Toleranz von 0,05 Euro je
%     % Megawatt und Stunde senken, ohne dass die Vollverdraengung bricht, doch
%     % die Lage des Minimums haengt an der Reihenfolge des Abstiegs. Gemessen im
%     % Modellrepository, MODELL.md Abschnitt 11.3, liegt die umgekehrte
%     % Reihenfolge am 11.02.2025 um 5,5 Prozent und am 26.08.2025 um 17,7
%     % Prozent unter der Standardreihenfolge, wobei beide Vektoren die Zusage
%     % tragen. Der ausgewiesene Wert ist damit die vorsichtigere Seite einer
%     % gemessenen Spanne und nicht der Preis. Die Zahlen gehoeren in die
%     % Ergebnisse, der Stichpunkt traegt deshalb keine.
%     % NICHT ALS FEHLER FORMULIEREN, sondern als Eigenschaft des Verfahrens.
%     \item Abhängigkeit des ausgewiesenen Preisvektors vom Suchverfahren, nämlich ein koordinatenweises und kein globales Minimum, dessen Lage von der Reihenfolge des Abstiegs abhängt
%     % F7 [daten] 69 Prozent der Jahressumme der RMS-Werte stammen aus
%     % dem groessten Prozent der Viertelstunden. Eine Groesse mit dieser
%     % Konzentration als Volatilitaetsmass fuer eine Verguetung zu
%     % verwenden, ist begruendungsbeduerftig.
%     \item Eigenschaft des Volatilitätsmaßes der Festlegung als Maß der Verteilungsränder
%     % F8 [eigen] Fuer eine Anlage mit rund 2,25 Stunden Entladedauer ist
%     % die Kopplung der Ladezustandsbilanz bindend, die unabhaengige
%     % Bewertung einzelner Zeitscheiben ignoriert sie.
%     \item Unabhängige Bewertung einzelner Zeitscheiben im Referenzverfahren gegenüber der bindenden Kopplung des Ladezustands
% \end{itemize}
%
%% ---------------------------------------------------------------
% \section{Handlungsempfehlungen}
% \label{sec:recommendations}
%% [Ziel: 1 Seite]
%%
%% Alle vier Empfehlungen sind [eigen] und aus den belegten Punkten
%% abgeleitet. Sie folgen jeweils aus einem Befund der vorangehenden
%% Abschnitte und sind ohne diesen nicht zu formulieren.
%
% \begin{itemize}
%     % G1, aus D3 und B6
%     \item Vergütung der Vorhaltung anstelle der Arbeit
%     % G2, aus D4
%     \item Nachweisverfahren für die Verfügbarkeit als konstitutiver Bestandteil des Mechanismus
%     % G3, aus C1
%     \item Behandlung der Doppelbelegung des Ladezustandsbands gegenüber der Vorhaltung von Regelleistung
%     % G4, aus D5
%     \item Umgang mit lokaler Marktmacht an Netzknoten mit geringer Anbieterzahl
% \end{itemize}
